\documentclass[12pt,oneside]{book}

% ============================================================
% METADATOS
% ============================================================
\newcommand{\universidad}{Universidad de Buenos Aires (UBA)}
\newcommand{\facultad}{Facultad de Derecho}
\newcommand{\materia}{Derecho Civil --- Introducción, Parte General}
\newcommand{\materiaabrev}{Derecho Civil}
\newcommand{\titulo}{Los efectos de la ley con relación al tiempo}
\newcommand{\autor}{Isaac Edgar Camacho Ocampo}
\newcommand{\anio}{2026}

% ============================================================
% PAQUETES
% ============================================================
\usepackage[utf8]{inputenc}
\usepackage[T1]{fontenc}
\usepackage[spanish,es-nodecimaldot]{babel}

\usepackage[
    top=2.5cm,
    bottom=2.5cm,
    left=3cm,
    right=2.5cm,
    headheight=15pt
]{geometry}

\usepackage{amsmath}
\usepackage{amssymb}
\usepackage{mathtools}

\usepackage{graphicx}
\usepackage{float}
\usepackage{caption}
\usepackage{subcaption}

\usepackage{booktabs}
\usepackage{array}
\usepackage{longtable}
\usepackage{tabularx}

\usepackage{enumitem}

\usepackage{xcolor}
\usepackage[most]{tcolorbox}

\usepackage{fancyhdr}
\usepackage{ragged2e}

\graphicspath{
    {./img/}
    {./graficos/}
    {./figuras/}
}

% ============================================================
% CONFIGURACIÓN
% ============================================================
\setcounter{tocdepth}{2}

\setlength{\parindent}{0pt}
\setlength{\parskip}{0.5em}

\setlist[itemize]{
    topsep=0.3em,
    itemsep=0.2em
}
\setlist[enumerate]{
    topsep=0.3em,
    itemsep=0.2em
}

\clubpenalty=10000
\widowpenalty=10000

\pagestyle{fancy}
\fancyhf{}
\fancyhead[L]{\nouppercase{\leftmark}}
\fancyhead[R]{\materiaabrev}
\fancyfoot[C]{\thepage}
\renewcommand{\headrulewidth}{0.4pt}
\renewcommand{\footrulewidth}{0pt}

\fancypagestyle{plain}{
    \fancyhf{}
    \fancyfoot[C]{\thepage}
    \renewcommand{\headrulewidth}{0pt}
}

% ============================================================
% COMANDOS
% ============================================================
\newcommand{\abs}[1]{\left|#1\right|}
\newcommand{\norm}[1]{\left\lVert#1\right\rVert}
\newcommand{\vect}[1]{\mathbf{#1}}
\newcommand{\deriv}[2]{\frac{d#1}{d#2}}
\newcommand{\pderiv}[2]{\frac{\partial#1}{\partial#2}}

% ============================================================
% ENTORNOS / CAJAS ACADÉMICAS
% ============================================================
\newtcolorbox{definition}[1]{
    enhanced,
    breakable,
    title={Definición: #1},
    fonttitle=\bfseries,
    colback=blue!3,
    colframe=blue!50!black,
    boxrule=0.8pt,
    arc=2mm
}

\newtcolorbox{important}{
    enhanced,
    breakable,
    title={Importante},
    fonttitle=\bfseries,
    colback=yellow!8,
    colframe=orange!70!black,
    boxrule=0.8pt,
    arc=2mm
}

\newtcolorbox{example}[1]{
    enhanced,
    breakable,
    title={Ejemplo: #1},
    fonttitle=\bfseries,
    colback=green!4,
    colframe=green!40!black,
    boxrule=0.8pt,
    arc=2mm
}

\newtcolorbox{formula}[1]{
    enhanced,
    breakable,
    title={Fórmula / Regla: #1},
    fonttitle=\bfseries,
    colback=purple!4,
    colframe=purple!50!black,
    boxrule=0.8pt,
    arc=2mm
}

\newtcolorbox{exercise}[1]{
    enhanced,
    breakable,
    title={Ejercicio: #1},
    fonttitle=\bfseries,
    colback=gray!5,
    colframe=gray!60!black,
    boxrule=0.8pt,
    arc=2mm
}

\newtcolorbox{note}{
    enhanced,
    breakable,
    title={Nota},
    fonttitle=\bfseries,
    colback=white,
    colframe=gray!50,
    boxrule=0.6pt,
    arc=2mm
}

% ============================================================
% HIPERVÍNCULOS Y METADATOS PDF
% ============================================================
\usepackage[
    colorlinks=true,
    linkcolor=blue!50!black,
    citecolor=green!40!black,
    urlcolor=blue!60!black,
    pdftitle={\titulo},
    pdfauthor={\autor},
    pdfsubject={Apuntes universitarios},
    bookmarks=true
]{hyperref}

% ============================================================
% DOCUMENTO
% ============================================================
\begin{document}

% ---------- PORTADA ----------
\begin{titlepage}
\thispagestyle{empty}

\begin{center}

\vspace*{1cm}

{\large \textsc{\universidad}}

\vspace{0.2cm}

{\large \textsc{\facultad}}

\vspace{3cm}

{\Huge \textbf{\materiaabrev}}

\vspace{1cm}

{\LARGE \titulo}

\vspace{0.8cm}

\textit{Comprender cómo el derecho organiza el paso del tiempo es comprender una de las claves de toda relación jurídica.}

\vspace{1.2cm}

\rule{0.70\textwidth}{0.5pt}

\vspace{0.5cm}

{\large Apuntes de estudio}

\vspace{0.5cm}

\rule{0.70\textwidth}{0.5pt}

\vfill

{\large \autor}

\vspace{0.3cm}

{\large \anio}

\end{center}

\vfill

\begin{flushleft}
\textit{Este es un modesto aporte a los estudiantes de ``\materiaabrev'' no pretende sustituir el material oficial de la materia.}
\end{flushleft}

\end{titlepage}

% ---------- ÍNDICE ----------
\tableofcontents

% ============================================================
% CONTENIDO
% ============================================================
\chapter{Los efectos de la ley con relación al tiempo}

Esta unidad aborda una cuestión central del derecho civil argentino: qué sucede cuando una ley cambia mientras una relación o situación jurídica se extiende en el tiempo. El problema recorre tres momentos históricos: el sistema del Código Civil de Vélez Sarsfield (1871), la reforma introducida por la Ley 17.711 (1968) y el régimen vigente en el artículo 7 del Código Civil y Comercial de la Nación.

Las cuatro unidades temáticas que componen esta clase forman, en conjunto, una sola historia que avanza en el tiempo, del mismo modo que el problema que describen: primero se plantea el conflicto; luego se examina cómo lo resolvía el Código originario, atendiendo a si el derecho ya estaba \textbf{adquirido} o era apenas una \textbf{expectativa}; a continuación se explica por qué esa solución resultó insuficiente y fue reemplazada por el criterio de los hechos \textbf{cumplidos}; y finalmente se muestra que el Código vigente conserva esa misma lógica, con un agregado orientado a la protección del consumidor.

\section{Planteamiento del problema}

La pregunta que estructura toda la unidad puede formularse así: ¿qué sucede si hay cambios en la legislación aplicables a las situaciones y relaciones jurídicas que se extienden en el tiempo? ¿Puede el legislador dictar normas retroactivas? ¿Con qué límites puede hacerlo?

\begin{example}{El préstamo en dólares}
Una persona presta a otra diez mil dólares a una tasa de interés del cinco por ciento anual, mediante un contrato de mutuo regido por el Código Civil. Mientras la relación está en curso, el Código Civil cambia, y la nueva ley dispone que quien prestó dólares debe recibir la devolución en pesos. Se genera así un conflicto entre las leyes a lo largo del tiempo: la ley anterior decía una cosa y la nueva ley dice otra, mientras que la relación jurídica no nació y concluyó bajo una misma ley, sino que se desarrolló en el tiempo y, durante ese desarrollo, se produjo un cambio normativo.
\end{example}

\subsection{Tres modos de aplicación de la ley en el tiempo}

A partir de este ejemplo pueden distinguirse tres situaciones posibles frente a un cambio legislativo:

\begin{enumerate}
    \item \textbf{Aplicación inmediata}: la nueva ley rige las consecuencias no cumplidas de relaciones jurídicas en curso.
    \item \textbf{Aplicación retroactiva}: la nueva ley pretende regir situaciones jurídicas ya concluidas bajo la ley anterior.
    \item \textbf{Aplicación para el futuro}: la nueva ley rige únicamente las relaciones jurídicas que nazcan a partir de su vigencia.
\end{enumerate}

\begin{example}{Aplicación retroactiva sobre un préstamo ya cumplido}
Si la misma persona hubiera prestado el dinero en enero de 2015 y se le hubiera devuelto en julio de 2015, y el nuevo Código --- vigente desde agosto de 2015 --- pretendiera aplicarse a lo ya cumplido, se estaría frente a una aplicación retroactiva de la ley. La tercera hipótesis, la aplicación para el futuro, es la que rige únicamente a partir de la entrada en vigencia del nuevo código, hacia adelante.
\end{example}

\begin{example}{El divorcio y los alimentos}
Bajo el Código Civil anterior, un cónyuge declarado culpable en un divorcio podía ser condenado a pagar alimentos a su cónyuge inocente. El Código Civil y Comercial, sancionado en 2015, eliminó la figura de la culpa en el divorcio y dispuso que ya no correspondía el pago de alimentos por esa causa. Frente a los reclamos de quienes debían seguir pagando alimentos bajo la ley anterior, los jueces entendieron que la nueva ley se aplicaba de inmediato a la situación jurídica en curso, de modo que el obligado dejaba de pagar alimentos desde la entrada en vigencia de la nueva norma. Queda así planteado el conflicto entre la ley anterior --- que disponía el pago de alimentos al cónyuge inocente --- y la nueva ley que lo eliminó, frente a una relación jurídica (el matrimonio y su disolución) desarrollada bajo la vigencia de la ley anterior.
\end{example}

\section{El sistema del Código Civil de Vélez Sarsfield (1871)}

El Código Civil originario, sancionado como Ley 340 en 1871, regulaba esta materia en su \textbf{artículo 3}, inspirado en el Código Civil francés.

\begin{formula}{Artículo 3 --- Código Civil de Vélez Sarsfield (texto original)}
``La ley ha de disponer para lo futuro; no tiene efecto retroactivo, ni puede alterar los derechos adquiridos.''
\end{formula}

Este artículo consagra el \textbf{principio de irretroactividad de la ley}: en principio, una ley que se sanciona rige el presente y el futuro, y no alcanza a lo que ya ha ocurrido. Todo el sistema del Código de Vélez estaba organizado en torno a la noción de \textbf{derecho adquirido}.

\subsection{El derecho adquirido}

\begin{definition}{Derecho adquirido}
Un derecho se considera adquirido cuando se cumplieron todos los recaudos que la ley de fondo dispone para que ese derecho ingrese al patrimonio de la persona, de modo que pueda considerarse su titularidad. El derecho adquirido se distingue de otras dos figuras: el derecho en expectativa (artículo 4044) y la mera facultad (artículo 4045).
\end{definition}

\subsection{El derecho en expectativa (artículo 4044)}

\begin{formula}{Artículo 4044 --- Código Civil de Vélez Sarsfield}
La nueva ley que entra en vigencia no se considera retroactiva si priva a los particulares de meros derechos en expectativa.
\end{formula}

\begin{definition}{Derecho en expectativa}
Es aquel derecho respecto del cual todavía no se han configurado los presupuestos necesarios para que la adquisición ocurra. Por esa razón, no se considera protegido frente a una nueva ley.
\end{definition}

\begin{example}{La expectativa de heredar}
Si una persona tiene la expectativa de heredar a un familiar y una nueva ley cambia el régimen de modo que ya no puede heredar, no puede sostener que esa ley sea retroactiva, porque el familiar todavía no ha muerto y, en consecuencia, todavía no se produjo la adquisición del patrimonio. Distinto sería el caso en que el familiar ya hubiera muerto y la nueva ley dispusiera que, de todas las personas fallecidas en los últimos dos años, solo heredan determinados familiares, excluyendo a quien reclama: en ese caso sí habría retroactividad sobre un derecho ya adquirido, porque ya se habían cumplido todos los recaudos para que el derecho sucesorio ingresara al patrimonio del heredero.
\end{example}

\subsection{La mera facultad (artículo 4045)}

\begin{formula}{Artículo 4045 --- Código Civil de Vélez Sarsfield}
Si la ley anterior otorgaba una facultad de reclamar algo y esa facultad no fue ejercida, la nueva ley que priva de esa facultad no es retroactiva.
\end{formula}

\begin{example}{La facultad no ejercida}
Si la ley anterior otorgaba la facultad de reclamar algo y esa facultad no fue ejercida en su momento, y la nueva ley priva de esa facultad, no puede sostenerse que la nueva ley sea retroactiva, porque quien tenía la facultad no la ejerció y, por lo tanto, se lo priva de algo que en rigor no llegó a tener.
\end{example}

\subsection{Las excepciones del sistema de Vélez}

El sistema se completaba con dos artículos adicionales, considerados tradicionalmente excepciones al principio de irretroactividad: los artículos 4 y 5.

\paragraph{a) Las leyes interpretativas (artículo 4) --- una falsa excepción.}

\begin{formula}{Artículo 4 --- Código Civil de Vélez Sarsfield}
Las leyes que tengan por objeto aclarar o interpretar otras leyes no tienen efecto respecto de los casos ya juzgados.
\end{formula}

\begin{example}{Ley interpretativa sobre la tasa de interés}
Si una ley dispone que la tasa de interés no puede superar un monto razonable, y al año siguiente otra ley aclara que ese monto razonable no puede superar el cinco por ciento, esta segunda ley está interpretando a la primera. La ley interpretativa se aplica a los casos en curso, pero no a los casos ya juzgados: existe allí un límite a lo que podría llamarse la retroactividad de la ley interpretativa.
\end{example}

\begin{note}
La doctrina considera que esta es una \textbf{falsa excepción}, porque, al ser interpretativa, en realidad se sigue aplicando la ley originaria; lo único que cambia es la interpretación que se le da a esa ley.
\end{note}

\paragraph{b) Las leyes de orden público (artículo 5) --- la verdadera excepción.}

\begin{formula}{Artículo 5 --- Código Civil de Vélez Sarsfield}
Ninguna persona puede tener derechos irrevocablemente adquiridos contra una ley de orden público.
\end{formula}

\begin{important}
El artículo 5 constituye una excepción genuina al principio de irretroactividad: si el sistema básico es la irretroactividad, una nueva ley de orden público --- es decir, que expresa principios fundamentales relativos a la subsistencia de la organización social --- puede aplicarse retroactivamente. Esta posibilidad generaba una considerable inseguridad jurídica, por lo que la Corte Suprema estableció que la ley de orden público no podía afectar derechos como la propiedad, garantizados por la Constitución.
\end{important}

\subsection{Síntesis del sistema de Vélez Sarsfield}

\begin{table}[H]
\centering
\begin{tabularx}{\textwidth}{>{\bfseries\raggedright\arraybackslash}p{0.18\textwidth} X}
\toprule
\textbf{Artículo} & \textbf{Contenido} \\
\midrule
Art. 3 & Principio de irretroactividad de la ley y protección de los derechos adquiridos. \\
Art. 4044 & El derecho en expectativa no está protegido frente a la nueva ley. \\
Art. 4045 & La mera facultad no ejercida no está protegida frente a la nueva ley. \\
Art. 4 & Las leyes interpretativas se aplican a casos en curso, no a casos juzgados (falsa excepción). \\
Art. 5 & Las leyes de orden público pueden afectar derechos adquiridos (excepción real). \\
\bottomrule
\end{tabularx}
\caption{Síntesis del sistema del Código Civil de Vélez Sarsfield}
\end{table}

\subsection{Crítica al sistema de Vélez}

Este sistema fue objeto de una crítica que se volvió clásica, siguiendo al autor francés Roubier, según la cual la teoría de los derechos adquiridos no servía para explicar cuándo una ley era retroactiva, especialmente en los casos de situaciones o relaciones jurídicas que se prolongan en el tiempo.

\begin{example}{El préstamo en cuotas}
Si un préstamo se pactó en varias cuotas, la primera ya cumplida y la segunda pendiente, y cambia la ley disponiendo que en lugar de dólares se puede entregar pesos, se genera un conflicto entre el deudor --- que pretende pagar en pesos --- y el acreedor --- que pretende cobrar en dólares, aplicando la ley anterior. La idea de derecho adquirido no permite resolver este conflicto, porque el acreedor sostendrá que el derecho ya está adquirido desde el primer momento, mientras que el deudor sostendrá que ese derecho no está adquirido sino que constituye una mera expectativa.
\end{example}

\section{La reforma de la Ley 17.711 (1968)}

La Ley 17.711, sancionada en 1968, constituye el punto de inflexión decisivo en esta materia, dado que es, en lo sustancial, el sistema que continúa rigiendo en la Argentina luego de que la reforma del Código Civil y Comercial mantuviera este texto en su parte central. La reforma modificó el \textbf{artículo 3} y derogó los artículos 4, 5, 4044 y 4045.

\begin{important}
La Ley 17.711 abandona la noción de derechos adquiridos y adopta la \textbf{teoría de los hechos cumplidos, o del consumo jurídico}, como criterio para determinar la retroactividad de la ley.
\end{important}

\begin{definition}{Teoría de los hechos cumplidos o del consumo jurídico}
Un hecho jurídico es una contingencia susceptible de producir consecuencias jurídicas. Cuando un hecho está totalmente cumplido --- es decir, cuando se ha consumido, cuando ha agotado la virtualidad de producir consecuencias jurídicas porque ya no puede producir más ---, se dice que se ha producido el consumo jurídico. Este criterio permite comprender que existe un elemento dinámico en las relaciones y situaciones jurídicas: los hechos pueden agotar todas sus consecuencias jurídicas en un instante, o pueden cumplirse a lo largo del tiempo.
\end{definition}

\begin{example}{Bicicleta versus venta en cuotas}
La compra de un bien mueble, como una bicicleta, en la que se paga y se entrega la cosa en un mismo acto, agota sus consecuencias jurídicas de manera instantánea. En cambio, una venta pactada en cuotas se cumple a lo largo del tiempo: la relación jurídica se consume recién cuando se paga la última cuota.
\end{example}

La teoría de los hechos cumplidos explica, mejor que la del derecho adquirido, qué sucede con la ley cuando existen relaciones y situaciones jurídicas que se extienden a lo largo del tiempo.

\begin{formula}{Artículo 3 --- Código Civil, texto según la reforma de la Ley 17.711}
``A partir de su entrada en vigencia, las leyes se aplican a las consecuencias de las relaciones y situaciones jurídicas existentes.''

``Las leyes no tienen efecto retroactivo, sean o no de orden público, salvo disposición en contrario. La retroactividad establecida por la ley en ningún caso podrá afectar derechos amparados por garantías constitucionales.''

``A los contratos en curso de ejecución no son aplicables las nuevas leyes supletorias.''
\end{formula}

Este artículo incorpora tres principios fundamentales, que reemplazan a los cinco artículos del sistema originario.

\subsection{Primer principio: efecto inmediato de la nueva ley imperativa}

\begin{definition}{Efecto inmediato}
A partir de su entrada en vigencia, la nueva ley se aplica de inmediato a las \textbf{consecuencias} de las relaciones y situaciones jurídicas existentes, consecuencias que se extienden en el tiempo. Este principio se refiere a las \textbf{leyes imperativas}, es decir, aquellas que se imponen a los particulares y que estos no pueden dejar de lado.
\end{definition}

\begin{example}{Cuota pendiente de un préstamo}
Si se prestó dinero y la primera cuota fue pagada mientras la segunda no, esta última constituye una consecuencia no cumplida de la relación jurídica; por lo tanto, la nueva ley se aplica a esa consecuencia pendiente.
\end{example}

\subsection{Segundo principio: irretroactividad de la ley}

Este principio se desagrega en cuatro componentes:

\begin{enumerate}
    \item Las leyes no tienen efecto retroactivo: la ley se aplica para el futuro, no hacia el pasado.
    \item Deja de ser relevante si la ley es o no de orden público, en coincidencia con la derogación del artículo 5.
    \item La propia ley puede disponer su retroactividad (``excepto disposición en contrario'').
    \item La retroactividad que establezca la ley no puede afectar derechos amparados por garantías constitucionales.
\end{enumerate}

Respecto del segundo componente, la irretroactividad y el orden público dejan de confundirse: se trata siempre de leyes imperativas, pero la imperatividad puede o no coincidir con el orden público, tanto en el Código Civil como en el Código Civil y Comercial.

\begin{example}{Retroactividad expresa: la filiación por técnicas de reproducción humana asistida}
El Código Civil y Comercial, al sancionarse, aplicó retroactivamente las normas sobre filiación derivada de la fecundación artificial a personas nacidas con anterioridad al Código, con el fin de solucionar el problema filiatorio de esas personas. Se trata de un caso de retroactividad establecida expresamente por la ley.
\end{example}

\begin{important}
El legislador tiene la potestad de ordenar hacia el pasado, pero esa potestad tiene un límite: la retroactividad que establezca la ley no puede afectar derechos amparados por garantías constitucionales. Si una ley dispusiera que todos los préstamos en dólares de los últimos dos años deben transformarse a pesos, ello afectaría un derecho de propiedad ya adquirido.
\end{important}

\subsection{Tercer principio: efecto diferido de la nueva ley supletoria}

\begin{formula}{Artículo 3, último párrafo --- texto original de la reforma de la Ley 17.711}
``Las nuevas leyes supletorias no son aplicables a los contratos en curso de ejecución.''
\end{formula}

\begin{definition}{Ley imperativa y ley supletoria}
La \textbf{ley supletoria} es aquella que las partes de un contrato pueden dejar de lado mediante pacto en contrario; si nada dicen, se aplica lo que la ley dispone. La \textbf{ley imperativa}, en cambio, se impone a las partes, que carecen de la libertad de apartarse de ella.
\end{definition}

\begin{example}{La sublocación}
En el Código Civil y Comercial, si las partes de una locación no dicen nada respecto de la posibilidad de sublocar, se entiende que sí pueden hacerlo; pero las partes pueden pactar algo distinto. Si el Código cambia esa regla supletoria, el cambio no rige la cláusula que las partes ya integraron a su contrato, porque son libres de disponer un régimen distinto mientras la norma sea supletoria.
\end{example}

\begin{example}{El plazo mínimo de la locación de vivienda}
El contrato de alquiler tiene un plazo mínimo que las partes no pueden dejar de lado: no pueden pactar apartarse de ese mínimo, porque la locación destinada a vivienda tiene un plazo mínimo de dos años. Esa imperatividad de la norma es distinta del carácter de la ley supletoria.
\end{example}

En consecuencia, una nueva ley supletoria tiene efecto diferido: no se aplica de inmediato, sino recién a partir de las nuevas relaciones jurídicas que se celebren; no se aplica a un contrato que ya está en curso de ejecución.

\section{El Código Civil y Comercial vigente (artículo 7)}

El Código Civil y Comercial, vigente desde 2015, mantuvo en lo sustancial la redacción de la reforma de 1968, ahora en su \textbf{artículo 7}.

\begin{formula}{Artículo 7 --- Código Civil y Comercial de la Nación}
``A partir de su entrada en vigencia, las leyes se aplican a las consecuencias de las relaciones y situaciones jurídicas existentes.''

``Las leyes no tienen efecto retroactivo, sean o no de orden público, excepto disposición en contrario. La retroactividad establecida por la ley no puede afectar derechos amparados por garantías constitucionales.''

``Las nuevas leyes supletorias no son aplicables a los contratos en curso de ejecución, con excepción de las normas más favorables al consumidor en las relaciones de consumo.''
\end{formula}

Los dos primeros principios no presentan cambios respecto del artículo 3 anterior: se mantiene el principio de efecto inmediato y el principio de irretroactividad, sin que la ley pueda afectar mediante retroactividad los derechos amparados por garantías constitucionales.

\subsection{La novedad: las normas más favorables al consumidor}

El cambio introducido por el artículo 7 se encuentra en el tercer principio.

\begin{important}
Las nuevas leyes supletorias no son aplicables a los contratos en curso de ejecución, \textbf{con excepción de las normas más favorables al consumidor en las relaciones de consumo}. Es decir que, si el contrato en curso de ejecución es un contrato de consumo, la nueva norma más favorable al consumidor se aplica de inmediato.
\end{important}

\begin{note}
Se señaló como observación crítica que este agregado transforma, en los hechos, la norma del consumidor en una norma imperativa, dado que se aplica de inmediato, y las normas de protección al consumidor suelen ser imperativas. Desde este punto de vista, el agregado resultaría un tanto inconsistente técnicamente, porque se formula como una excepción referida a normas supletorias, cuando en rigor la norma favorable al consumidor en una relación de consumo suele ser de carácter imperativo. Con todo, el cambio se explica por la incorporación del derecho del consumidor al nuevo Código Civil y Comercial, ausente en el código anterior.
\end{note}

\section{Cuadro comparativo de los tres sistemas}

\begin{table}[H]
\centering
\begin{tabularx}{\textwidth}{>{\bfseries\raggedright\arraybackslash}p{0.22\textwidth} X X}
\toprule
\textbf{Aspecto} & \textbf{Vélez (1871) --- arts. 3, 4, 5, 4044, 4045} & \textbf{17.711 (1968) / CCyC (2015) --- arts. 3 y 7} \\
\midrule
Criterio rector & Derechos adquiridos & Hechos cumplidos / consumo jurídico \\
Irretroactividad & Principio general (art. 3) & Principio general; admite retroactividad expresa con límite constitucional \\
Orden público & Excepción real a la irretroactividad (art. 5) & Irrelevante a estos fines: no se confunde con la imperatividad \\
Leyes interpretativas & Falsa excepción (art. 4) & No reguladas como categoría aparte \\
Leyes supletorias y contratos en curso & No reguladas expresamente & Efecto diferido; excepción a favor del consumidor (solo CCyC, art. 7) \\
\bottomrule
\end{tabularx}
\caption{Comparación de los tres sistemas normativos}
\end{table}

\begin{note}
\textbf{Recapitulación final expuesta por el docente:}
\begin{itemize}
    \item El problema: qué ocurre con la ley cuando una relación o situación jurídica se extiende en el tiempo y sobreviene un cambio normativo.
    \item El sistema de Vélez Sarsfield (1871): basado en la noción de derecho adquirido, distinguida de la expectativa y de la mera facultad, con dos excepciones (leyes interpretativas y leyes de orden público).
    \item La reforma de la Ley 17.711 (1968): abandona los derechos adquiridos y adopta la teoría de los hechos cumplidos o del consumo jurídico, articulada en tres principios --- efecto inmediato, irretroactividad y efecto diferido de la ley supletoria.
    \item El Código Civil y Comercial vigente: mantiene la redacción de 1968 en el artículo 7, y agrega como única novedad la aplicación inmediata de las normas más favorables al consumidor en los contratos de consumo en curso.
\end{itemize}
\end{note}

\section*{Cuadro Cornell}
\addcontentsline{toc}{section}{Cuadro Cornell}

El siguiente cuadro sintetiza, en formato Cornell, los conceptos clave abordados a lo largo de la unidad.

\begin{longtable}{>{\bfseries\raggedright\arraybackslash}p{0.30\textwidth} p{0.60\textwidth}}
\toprule
\textbf{Preguntas / palabras clave} & \textbf{Notas / respuestas} \\
\midrule
\endfirsthead
\toprule
\textbf{Preguntas / palabras clave} & \textbf{Notas / respuestas} \\
\midrule
\endhead
\bottomrule
\endfoot
\bottomrule
\endlastfoot

¿Cuál es el problema central de la unidad? &
Determinar qué sucede cuando una ley cambia mientras una situación o relación jurídica se está desarrollando en el tiempo: si la nueva ley se aplica de inmediato, retroactivamente o solo hacia el futuro. \\
\midrule

¿Qué es un derecho adquirido (sistema de Vélez, art. 3)? &
Aquel respecto del cual ya se cumplieron todos los recaudos que la ley exige para que ingrese al patrimonio de la persona. No puede ser alterado por una ley posterior, salvo excepciones. \\
\midrule

¿Qué es un derecho en expectativa (art. 4044)? &
Un derecho cuyos presupuestos de adquisición todavía no se cumplieron (por ejemplo, la expectativa de heredar a un familiar que aún vive). La nueva ley puede privarlo sin que ello implique retroactividad. \\
\midrule

¿Qué es una mera facultad (art. 4045)? &
Una facultad otorgada por la ley anterior que no fue ejercida en su momento. La nueva ley puede suprimirla sin que ello sea retroactividad. \\
\midrule

¿Por qué las leyes interpretativas (art. 4) son una falsa excepción? &
Porque, al interpretar la ley anterior, en rigor se sigue aplicando esa misma ley originaria, solo que a la luz de la interpretación dada. Se aplican a casos en curso, no a casos ya juzgados. \\
\midrule

¿Por qué las leyes de orden público (art. 5) eran la verdadera excepción? &
Porque, al expresar principios fundamentales de la organización social, podían aplicarse retroactivamente y afectar derechos adquiridos, generando inseguridad jurídica (limitada luego por la Corte Suprema en defensa de la propiedad). \\
\midrule

¿Por qué se reformó el sistema en 1968 (Ley 17.711)? &
Porque la teoría de los derechos adquiridos no explicaba adecuadamente los casos de relaciones jurídicas que se desarrollan en el tiempo (por ejemplo, un préstamo en cuotas), donde no estaba claro si el derecho ya estaba adquirido o era solo una expectativa. \\
\midrule

¿Qué es la teoría de los hechos cumplidos o del consumo jurídico? &
El criterio que reemplaza al derecho adquirido: distingue si un hecho jurídico ya agotó (consumió) todas sus consecuencias, o si todavía tiene consecuencias pendientes a las que se aplicará la nueva ley. \\
\midrule

¿Cuáles son los tres principios del artículo 3 (17.711) y del artículo 7 (CCyC)? &
1) Efecto inmediato de la nueva ley imperativa sobre las consecuencias no cumplidas. 2) Irretroactividad de la ley, salvo disposición en contrario y sin afectar garantías constitucionales. 3) Efecto diferido de la nueva ley supletoria, que no se aplica a contratos en curso de ejecución. \\
\midrule

¿Qué diferencia hay entre ley imperativa y ley supletoria? &
La ley imperativa se impone a las partes, que no pueden dejarla de lado (por ejemplo, el plazo mínimo de la locación de vivienda). La ley supletoria rige solo si las partes no pactaron algo distinto, y puede ser dejada de lado por ellas. \\
\midrule

¿Qué novedad introdujo el Código Civil y Comercial en el artículo 7? &
Que las normas más favorables al consumidor se aplican de inmediato a los contratos de consumo en curso de ejecución, aun cuando se trate de leyes supletorias --- lo que se señaló como una excepción técnicamente algo inconsistente. \\
\midrule

¿Puede la ley ser retroactiva expresamente? &
Sí, si la propia ley lo dispone (``excepto disposición en contrario''), como ocurrió con las normas de filiación por técnicas de reproducción humana asistida en el CCyC, aplicadas a personas ya nacidas. El límite siempre son los derechos amparados por garantías constitucionales. \\
\midrule

\multicolumn{2}{p{0.90\textwidth}}{
\textbf{Resumen:} La clase recorrió la evolución histórica de un mismo problema: cómo trata el derecho argentino los conflictos que surgen cuando una ley cambia mientras una relación jurídica está en desarrollo. El Código de Vélez Sarsfield (1871) resolvía esta cuestión con la teoría de los derechos adquiridos, distinguiendo entre derechos ya adquiridos, meras expectativas y meras facultades, y admitiendo dos excepciones al principio de irretroactividad: las leyes interpretativas (una excepción solo aparente) y las leyes de orden público (una excepción real, aunque generadora de inseguridad jurídica). Esta construcción resultó insuficiente para explicar los casos de relaciones jurídicas que se desenvuelven en el tiempo, como un préstamo pagadero en cuotas, lo que llevó a la reforma de la Ley 17.711 en 1968 a reemplazar el criterio de derecho adquirido por la teoría de los hechos cumplidos o del consumo jurídico, plasmada en tres principios todavía vigentes: el efecto inmediato de la nueva ley imperativa sobre las consecuencias no cumplidas, la irretroactividad de la ley --- salvo disposición expresa en contrario y siempre dentro del límite de las garantías constitucionales --- y el efecto diferido de la nueva ley supletoria frente a los contratos en curso de ejecución. El Código Civil y Comercial vigente desde 2015 conservó esta misma estructura en su artículo 7, sumando como única novedad relevante la aplicación inmediata de las normas más favorables al consumidor en las relaciones de consumo en curso, en sintonía con la incorporación del derecho del consumidor al nuevo ordenamiento.
} \\

\end{longtable}

\end{document}
